\documentclass[12pt]{article}
\usepackage[utf8]{inputenc}
\usepackage[T1]{fontenc}
\usepackage{hyperref}
\title{Central Lalo Language Resources for K-12 Curriculum Development}
\author{Compiled by Claude Sonnet 5.5 (claude-sonnet-5-5) for Andr\'as Kornai}
\date{}
\begin{document}
\maketitle

\section{Introduction}
Central Lalo (ISO 639-3 ywt, Glottocode xish1235), also called Xishanba Lalo, is a Western Yi language of west-central Yunnan, China. Wikipedia gives about 213,000 speakers, while Global Recordings Network gives 43,000. I could not tell which figure is current. UNESCO classes Lalo as vulnerable. The search found little in the way of classroom materials.

\section{Primary Educational Resources}
No primers, textbooks or school programmes were found in the search results.

\section{Linguistic and Reference Materials}
\begin{itemize}
\item Wikipedia, Lalo language: \url{https://en.wikipedia.org/wiki/Lalo_language}. It cites a grammar by Bj\"orverud (1998), based on fieldwork in Longjie Township, Weishan County.
\item Endangered Languages Project entry for Central Lalo: \url{https://www.endangeredlanguages.com/lang/5129}
\item Numeral dataset for Central Lalo (Chan numerals project): \url{https://lingweb.eva.mpg.de/channumerals/Lalo-Central.htm}
\end{itemize}

\section{Religious and Community Resources}
Global Recordings Network has a page for this language: \url{https://globalrecordings.net/en/language/ywt}. I did not confirm what recordings it holds.

\section{Digital and Audio Resources}
No audio corpora or software were found beyond the Global Recordings Network page above.

\section{Curriculum Development Context}
No evidence was found that the language is taught in schools. A curriculum would need new documentation work with speakers, and the 1998 grammar would be the main published reference.
\end{document}